\documentclass[11pt,a4paper]{article}
\usepackage[margin=25mm]{geometry}
\usepackage{amsmath,amssymb,amsthm,booktabs,hyperref}
\newtheorem{theorem}{Theorem}
\newtheorem{proposition}[theorem]{Proposition}
\title{The two finite chiral atlases: recovered results and source corrections}
\author{Human Compression project}\date{}
\begin{document}\maketitle
This manuscript gives a source-wide disposition of Hartley, Hubard and
Leemans, \emph{Two atlases of abstract chiral polytopes for small groups},
Ars Mathematica Contemporanea 5 (2012), 371--382. The first atlas and all
printed entries of the second atlas have complete accepted finite evidence.
Several printed numbers and one automorphism-group label require correction.
The larger domain stated in Section~4 also requires an explicit boundary:
the printed roster has 79 groups and omits one group satisfying that section's
literal condition. We do not invent a census for that omitted group.

The appended complete component proofs and the exact source ledger specify
what is acquired. Standard regular and chiral reconstruction theorems retain
the source's attribution to McMullen--Schulte and Schulte--Weiss. Exact finite
GAP group operations, the named group constructions and automorphism
classifications, and completeness of the earlier regular-cover atlas remain
stated inputs. No present website is asserted to have identical historical
bytes. This synthesis repeats no group census.

\paragraph{Definitions and equivalence.}
We use the source's ranked-poset definition, including the diamond axiom
and strong flag connectivity. Its Sections~2--3 group criteria are retained
with their complete subset intersection property. The accepted prefix-coordinate
manuscript supplies the interval/difference subgroup dictionary, reflection
criterion and complete stabilizer enumeration; it does not replace polytope
recognition by adjacent product orders alone. A colour-preserving isomorphism
of flag graphs is a polytope isomorphism. Reversing all colours gives duality,
which is not additionally identified in the atlas counts below.

In the first atlas, the two enantiomorphic markings are identified as one
unrooted chiral polytope. In the second atlas, a generating rotation marking
is taken modulo the full automorphism group of its abstract group; its mirror
is counted separately. The reflection criterion implies that this mirror
involution has no fixed chiral class. Thus division by two there forgets
orientation, not duality. Regular columns count full regular automorphism
groups, not rotation groups extracted from the chiral search.

\begin{theorem}[Theorem 3.1 and the source algorithm]
Let the actual flag graph of a quotient polytope be the right-coset graph
of $N\le W=\langle r_0,\ldots,r_{d-1}\rangle$, where the $r_i$ are its
specified flag adjacencies. Its colour automorphism group is $N_W(N)/N$.
The quotient is chiral precisely when $[W:N_W(N)]=2$ and all $r_i$ lie
outside $N_W(N)$. There is at most one index-two subgroup $H$ avoiding all
$r_i$. For that $H$, the possible chiral stabilizers are precisely
$N\lhd H$ not normal in $W$, subject to the actual polytope condition.
Conjugate stabilizers $N$ and $r_0Nr_0$ describe the same unrooted quotient.
\end{theorem}
\begin{proof}
A colour automorphism is determined by the image $Na$ of $N$ and must
send $Nw$ to $Naw$. This is well-defined exactly when $a$ normalizes $N$.
The resulting action is free, with $[W:N_W(N)]$ orbits. For an index-two
normalizer, adjacency by $r_i$ exchanges the two orbits exactly when $r_i$
is outside it. The subgroup avoiding every generator is the kernel of the
unique possible parity homomorphism $r_i\mapsto1\in C_2$. A subgroup normal
in that $H$ but not in $W$ has normalizer exactly $H$, so it has the two
stated conjugates. The complete coset and geometric interface proofs are in
\emph{Chiral quotients from one parity subgroup}.
\end{proof}
The source's core-free filter is optional: the kernel of the flag action is
$\operatorname{Core}_W(N)=N\cap r_0Nr_0$. It is not legitimate to discard
nontrivial cores when excluding arbitrary chiral quotients. Nor do we infer
that every monodromy quotient is a regular polytope. Minimal faithful-cover
language is used only with its specified flag-action interface.

\begin{theorem}[The first atlas and Table 1]
Relative to the acquired complete regular-cover atlas of orders at most
2000, excluding 1024 and 1536, there are exactly 48 rank-three and eight
rank-four chiral quotients, with distinct duals counted separately.
The four Table~1 markings and their reversals are the eight rank-four
examples. Their displayed cover and rotation-group structures are correct.
The vertex figures of the last two rows have types $\{6,18\}$ and $\{6,9\}$
and are among the verified rank-three chiral examples.
\end{theorem}
\begin{proof}
The accepted full source parser supplies 9,212 nondegenerate marked covers.
The acquired product argument excludes a bond of order two; ranks one and
two cannot be chiral. All normal stabilizers allowed by the preceding
criterion are considered, including nontrivial cores. Three exact completed
intervals partition the source roster. Their 8,986 normal candidates leave
756 after the necessary base-face intersection test.

Literal coset transport verifies 704,608 edges. Of the 756 candidates,
360 have full colour isomorphisms to the 56 accepted examples, covering all
of them; the other 396 have explicit negative witnesses, namely 350 diamond
failures and 46 disconnected flag residues. The examples themselves have
complete face-poset, chain, diamond and residue checks, and complete rooted
colour comparisons proving distinctness and chirality. This is the accepted
\emph{Identifying the quotient flag graphs} component, with no reliance on
an interrupted candidate or on abstract database identifiers for exclusions.

Table~1's six structural names were proved on the actual cover groups and
rotation quotients by literal generator-edge transports, normal subgroups,
complements and commuting factors. In particular the last two rows really
have the same named cover and rotation groups. The source convention is
$|D_9|=18$; its parentheses are preserved. The complete names and certificates
are in \emph{The six group descriptions in the small chiral atlas}.

For the remaining vertex-figure sentence, keep the base vertex in either
accepted rank-four flag graph and take its component using colours $1,2,3$.
Renumbering those colours gives the vertex-figure flag graph by the polytope
section definition. The saved components have respectively 216 and 108 flags.
Complete bijections to accepted rank-three entries 44 and 39 preserve all
648 and 324 coloured edges. Thus the sections are the asserted chiral
polytopes. Every vertex has flags from both global flag orbits (an adjacency
of colour one preserves the vertex and exchanges those orbits), so the chiral
automorphism group is vertex-transitive. The base-vertex calculation therefore
applies to all vertex figures of each example.
\end{proof}

\begin{theorem}[All printed almost-simple entries]
For the 79 distinct named group models in Tables~2--7, occurring in 88 rows,
the regular and oriented chiral counts are as follows. All rows agree rank
by rank except the regular rank-three entry for $\mathrm{PSL}_2(27)$,
which is 26 rather than 27, and the chiral rank-three entry for
$\mathrm{PSU}_4(2)\cong\mathrm{PSp}_4(3)$, which is 46 rather than 26.
\[
\begin{array}{c|rrrrrr|r}
\text{rank}&3&4&5&6&7&8&\text{total}\\\hline
\text{regular}&8319&810&61&16&3&1&9210\\
\text{oriented chiral}&15056&1890&28&0&0&0&16974
\end{array}
\]
The oriented chiral classes form 8487 mirror pairs.
\end{theorem}
\begin{proof}
The accepted regular application acquires 14 earlier complete censuses and
65 complete exact executions. Its full automorphism actions are verified
common actions on the named models; it does not use an inner action in place
of the full automorphism group. The acquired records bind their actual drivers,
models and outputs, with complete terminal rank lists; the older records'
runtime-file-trace limits remain explicit. A separate fixed-endpoint
calculation gives $312/12=26$ for the exceptional regular row.

The accepted chiral application acquires 72 complete whole searches, the
complete 25-branch $S_9$ union, and six complete fixed-first-class unions.
A fixed first list is proved to contain exactly one representative of each
relevant conjugacy class. Unions require its literal equality, not just
matching class indices. Each centralizer recursion is complete by the
acquired stabilizer theorem. Heredity and strict subgroup growth supply its
all-rank stopping boundary. Full automorphisms act freely on generating
markings, so division of the complete inner-orbit count by $|\mathrm{Out}(G)|$
is valid separately at each rank and type. An independent literal positive
certificate already distinguishes 46 chiral classes for the exceptional
unitary row, without using its census completeness or the outer-order division.

The accepted applications join these results to every printed occurrence
through literal named models. Nine repeated occurrences are removed only
after their group identifications; they contribute 95 regular and 50 chiral
counts to the raw family sums. Their removal gives the displayed totals.
No source count was a filter or premise of these enumerations.
\end{proof}

\paragraph{Family totals and the observations in Section 6.}
The exact family sums, retaining repeated occurrences between families, are
\[
\begin{array}{c|rrrrrr}
\text{Table}&2&3&4&5&6&7\\\hline
\text{regular}&1536&804&3903&1846&1202&14\\
\text{oriented chiral}&4780&1710&1882&3998&4242&412
\end{array}
\]
Besides the two row corrections, the printed aggregate arithmetic in
Table~3 (1702), Table~4 (2050), and Table~5 (1840 regular) is corrected by
these sums. Section~6's global 9205, 17114 and 8557 become respectively
9210, 16974 and 8487. Its PSL$_2$-family rank-at-least-four counts are
45 regular and 338 chiral, replacing 51 and 506. The comparative comments
about sporadic, unitary and Suzuki groups remain statements about this
finite roster and its stated convention. They do not compare all groups
or all genera. As in the source's aggregate, $\operatorname{Aut}(J_2)$ is
excluded: its order 1209600 lies beyond both stated cutoffs. The older
5265 regular objects identified duals; the historical
Conder genus counts, MAGMA timings and website descriptions remain attributed
context, not new calculations or interchangeable counting conventions.

\paragraph{Orders and named automorphism groups.}
The complete source application identifies every printed group and automorphism
label through acquired constructor/classification and literal transport or
outer-preimage routes. For $S\le G\le\operatorname{Aut}(S)$ with $S$ simple,
restriction identifies
$\operatorname{Aut}(G)\cong N_{\operatorname{Aut}(S)}(G)$: the unique simple
socle is characteristic, its centralizer is trivial, and a normalizing
automorphism extends by conjugation. The actual preimages and small outer
normalizers, not order equality alone, supply the names. The accepted $J_2:2$
join additionally has an actual outside involution complement.

The source's order 443510 for $M_{22}$ is corrected to 443520, and its order
7800 for $\mathrm{PSL}_2(25).2$ to 15600. The automorphism label
$P\Gamma L_3(3)$ for $\mathrm{PSU}_3(3)$ is corrected to $P\Gamma U_3(3)$,
of order 12096. In contrast $P\Gamma L_3(3)=PGL_3(3)$ has order
$(26\cdot24\cdot18)/2=5616$.

\paragraph{The historical-domain statement requires a correction.}
The abstract uses a 900000 cutoff and speaks of the group $\Gamma$ appearing
in the 1985 ATLAS. Section~4 instead gives a one-million cutoff and explicitly
requires only the simple socle $S$ to appear. Under that latter socle condition,
the accepted complete simple/outer roster and the explicitly retained CTblLib
transcription of the 1985 contents give 80 almost-simple isomorphism classes.
Exactly 79 have order below 900000; these are precisely the printed models.
The additional group is $\mathrm{Sp}_4(4)$, of order 979200, whose simple
socle has ordinary ATLAS page 44. Thus the printed roster does not exhaust
the larger Section~4 domain. This corrects the coverage claim, rather than
silently replacing its domain by the printed list.

The historical ATLAS transcription is an explicit premise, not a claimed
direct inspection of a book scan. We do not infer that every outer extension
had its own printed ATLAS table, nor what an unavailable historical program
executed. No polytope count, zero result, or new census for $\mathrm{Sp}_4(4)$
is asserted. The recovered theorem is the complete printed 79-group atlas
with the corrected entries, plus this exact larger-domain omission.

\paragraph{Attribution and disposition.}
The original constructions, examples and enumeration programme retain the
source authors' credit. The complete component proofs identify their borrowed
classification and software interfaces. The Human Compression project supplied
the finite applications and operational repairs, with individual contributions,
distinct reviews, the normalizer suggestion and advisory source readings
recorded in the component provenance. The only additional finite join in this synthesis is the two
vertex-figure maps. This is a corrected source disposition, submitted for a
separate whole-source review. It claims no new conjecture solution, Lean proof,
compression measurement or automatic publication status.
\end{document}
